\section*{Hoofdstuk \ref{ch:g4:light-straight-lines} --- Licht plant zich rechtlijnig voort}

\begin{solution}{exo:g4:light-straight-lines:1}
\omterm{def:g1:light-and-shadows:source}{Licht} plant zich langs rechte lijnen voort --- het buigt nooit om
hoeken. Waarnemingen: rechte zonnestralen door nevel of jaloezieën,
het scherpe blad van een zaklamp in stoffige \omterm{def:g2:air-around-us:air}{lucht}, het kaarsrechte
draadje van een laser, de bundels van een vuurtoren.
\end{solution}

\begin{solution}{exo:g4:light-straight-lines:2}
Een rechte lijn met een pijl die staat voor één draadje \omterm{def:g1:light-and-shadows:source}{licht}: ze
laat de weg zien die het \omterm{def:g1:light-and-shadows:source}{licht} neemt (kaarsrecht) en de richting waarin
het loopt.
\end{solution}

\begin{solution}{exo:g4:light-straight-lines:3}
Precies wanneer de vlam, alle drie de gaatjes en jouw oog op één
rechte lijn staan. De middelste kaart een vingerbreedte opzij schuiven
maakt het zicht kapot: \omterm{def:g1:light-and-shadows:source}{licht} maakt geen stapje opzij --- geen rechte
lijn, geen \omterm{def:g1:light-and-shadows:source}{licht}.
\end{solution}

\begin{solution}{exo:g4:light-straight-lines:4}
De stralen lopen recht van de lamp weg; wat het \omterm{def:g1:light-and-shadows:source}{licht} tegenhoudt stopt
er een deel van, en de stralen die net langs de rand scheren lopen
recht verder. De grens tussen bereikte en onbereikte grond wordt door
rechte lijnen getekend --- vandaar de scherpe rand.
\end{solution}

\begin{solution}{exo:g4:light-straight-lines:5}
Het \omterm{def:g1:light-and-shadows:source}{licht} verlaat de lamp, loopt recht naar de bladzijde, kaatst van
het papier terug en loopt recht van de bladzijde je oog in --- dat het
ontvangt. Twee rechte trajecten met één kaatsing ertussen.
\end{solution}

\begin{solution}{exo:g4:light-straight-lines:6}
\omterm{def:g3:sound-around-us:sound}{Geluid} spreidt zich uit als kringen in een vijver en spoelt om de
rand van de hoek heen tot in je oor. \omterm{def:g1:light-and-shadows:source}{Licht} houdt zich aan rechte
wegen, en geen rechte lijn verbindt het gezicht van de vriend door
baksteen heen met jouw oog --- dus komt de stem wel aan en het beeld
niet.
\end{solution}

\begin{solution}{exo:g4:light-straight-lines:7}
Alleen vanaf de plekken waar een rechte lijn door het sleutelgat de
lamp bereikt --- oog, sleutelgat en lamp op één lijn. Stap opzij en de
lijn breekt op de deur.
\end{solution}

\begin{solution}{exo:g4:light-straight-lines:8}
De flits: de reis van het \omterm{def:g1:light-and-shadows:source}{licht} telt als ogenblikkelijk, terwijl de
knal met \qty{340}{m} per \omterm{def:g2:measuring-time:second}{seconde} oversjokt --- drie stille \omterm{def:g2:measuring-time:second}{seconden}
per kilometer afstand.
\end{solution}

\begin{solution}{exo:g4:light-straight-lines:9}
Er verlaat niets het oog. Het eigen \omterm{def:g1:light-and-shadows:source}{licht} van de \omterm{def:g4:solar-system:star}{ster} legt de rechte
weg \emph{naar} jou af, en zien is dat opvangen --- het oog is een
net, geen lantaarn.
\end{solution}

\begin{solution}{exo:g4:light-straight-lines:10}
Zet het glanzende oppervlak in de deuropening, schuin (onder een halve
rechte hoek) ten opzichte van de gang. Een straal die langs de gang
loopt raakt de schuine spiegel en zijn weg vouwt: hij gaat de hoek om
en loopt recht door de deuropening de kamer in tot in jouw oog. Twee
van zulke vouwen, de een boven de ander, maken de periscoop van de
onderzeeër.
\end{solution}

\begin{solution}{exo:g4:light-straight-lines:11}
Je ziet het \omterm{def:g1:light-and-shadows:source}{licht} niet onderweg --- je ziet de ontelbare piepkleine
druppeltjes en stofjes die de bundel onderweg raakt, en die elk een
beetje \omterm{def:g1:light-and-shadows:source}{licht} naar je oog kaatsen. In volmaakt schone \omterm{def:g2:air-around-us:air}{lucht} is er
niets om te raken, en is de bundel van opzij onzichtbaar: alleen de
heldere plek waar hij landt verraadt hem.
\end{solution}
